%%%%%%%%%%  scan idx 179  |  folio 166  %%%%%%%%%%
\[
\operatorname{Hom}_{\mathcal A}(L,M)\longrightarrow\operatorname{Hom}_{\mathcal B}(\varphi(L),\varphi(M))
\]
\emph{est un
isomorphisme, et pour tout $n\ge1$ et tout $x\in\Ext^n_{\mathcal B}(\varphi(L),\varphi(M))$ il existe un épimorphisme $L'\to L$ de $\mathcal A$ qui efface} $x$.
\end{enumerate}
\item \emph{Les conditions suivantes sont équivalentes} :
\begin{enumerate}[label=(\roman*)]
\item \emph{pour tous $L,M\in\Ob\mathcal A$ et tout $n\in\mathbb Z$, la flèche canonique}
\[
\Ext^n_{\mathcal A}(L,M)\longrightarrow\Ext^n_{\mathcal B}(\varphi(L),\varphi(M))
\]
\emph{est un isomorphisme} ;
\item \emph{le foncteur canonique} (\emph{induit par} $\varphi$)
\[
\mathrm D^b(\mathcal A)\longrightarrow\mathrm D^b(\mathcal B)
\]
\emph{est pleinement fidèle et a pour image essentielle la sous-catégorie pleine de $\mathrm D^b(\mathcal B)$ formée des complexes $F$ tels que $H^i(F)$ soit dans l'image essentielle de $\varphi$ pour tout} $i$.
\end{enumerate}
\end{enumerate}

\emph{Preuve.} a) Pour l'implication (i)$\Rightarrow$(ii), il suffit de montrer ceci : soient $n\ge1$, $L\in\Ob\mathcal A$, et $y\in\Ext^n_{\mathcal A}(L,M)$, il existe alors un épimorphisme $L'\to L$ qui efface $y$. On peut, pour cela, utiliser l'interprétation de Yoneda ([18] III 3.2) : $y$ est la classe d'une suite exacte
\[
0\longrightarrow M\longrightarrow E_{n-1}\longrightarrow\ldots\longrightarrow E_0\longrightarrow L\longrightarrow0;
\]
l'épimorphisme $E_0\to L$ efface $y$ ([18] III 3.2.5 et 3.2.8). Prouvons (ii)$\Rightarrow$(i). Posons, pour abréger,
\[
\Ext^n_{\mathcal A}(.,M)=F^n,\qquad
\Ext^n_{\mathcal B}(\varphi(.),\varphi(M))=G^n.
\]
L'hypothèse d'effaçabilité signifie que l'on a, pour $n\ge1$, et $L\in\Ob\mathcal A$,
\[
G^n(L)=\varinjlim\operatorname{Coker}(G^{n-1}(L_0)\to G^{n-1}(L_1)),
\]
la limite inductive étant prise suivant l'ensemble des suites exactes 
